% ============================================================
% SM-3: K3 Spectral Theorem — Koide K = 2/3 from Colour Cage Base Graph
% 600-Cell Standard Model Emergence Series
%
% Version 6 — Layer A/B/C epistemic decomposition; robustness calculation
%
% v5→v6 changes (16 April 2026, ChatGPT referee review):
%   C1: New §3 "Epistemic Layer Structure" — Layer A (geometric input),
%       Layer B (imported open-system formalism), Layer C (mathematical
%       result).  Same architecture as SS-3 v1.4.
%   C2: P3 derivation revised: Caldeira–Leggett coupling, τ_relax
%       ordering, and full Gibbs equilibration explicitly labelled as
%       Layer B assumptions; open problem SS-4 cited as resolution path.
%   C3: Finite-temperature robustness calculation added with correct
%       O(ℏω₀/kT_P) ~ O(10⁻²⁰) scaling (correcting earlier e⁻¹⁰²⁰ claim).
%   C4: Scope table updated to reflect Layer A/B/C status.
%   C5: Acknowledgements updated (ChatGPT reviewer credited).
%   C6: Bibliography expanded (Caldeira–Leggett 1983, Breuer–Petruccione 2002).
%
% v4→v5 changes (Session D, 24 March 2026):
%   Issue 1: Added Postulate P3 (thermal state-counting equipartition)
%   Issue 2: Added Remark on why quarks do NOT satisfy Koide
%   Issue 3: Scope table made more prominent
%
% v5 harmonization (26 March 2026):
%   H1: Series ID SM-3 added to title; series name corrected
%   H2: Author line — Claude Sonnet (Anthropic); institution block added
%   H3: Bibliography added (Koide 1982, PDG 2024, SS-1, SM-1, SM-4)
%   H4: Acknowledgements section added
%   H5: Date updated to Version 5, 26 March 2026
%   H6: hyperref added
%   H7: OP-SM-7d registered with \openproblem environment
%   H8: siunitx and graphicx added
%   H9: Tagline removed from date block; email added
%   H10: GitHub URLs added to all CPP bibliography entries
% ============================================================

\documentclass[11pt]{article}
\usepackage{amsmath,amssymb,amsthm,geometry,booktabs,array,enumitem}
\usepackage{siunitx}
\usepackage{graphicx}
\usepackage{hyperref}
\usepackage[numbers,sort&compress]{natbib}
\hypersetup{colorlinks=true,linkcolor=blue,citecolor=blue,urlcolor=blue}
\geometry{margin=1in}

\newtheorem{theorem}{Theorem}[section]
\newtheorem{lemma}[theorem]{Lemma}
\newtheorem{corollary}[theorem]{Corollary}
\newtheorem{postulate}{Postulate}
\newtheorem{definition}{Definition}[section]
\newtheorem{remark}{Remark}[section]
\newtheorem{proposition}[theorem]{Proposition}
\newtheorem{openproblem}{Open Problem}

\newcommand{\phig}{\varphi}
\newcommand{\ZBW}{\mathrm{ZBW}}

\title{\textbf{SM-3: K3 Spectral Theorem and the Koide Formula}\\[6pt]
{\large 600-Cell Standard Model Emergence Series}\\[4pt]
{\normalsize Version 6, 16 April 2026}}

\author{Thomas Lee Abshier, ND \and Grok (xAI) \and Claude Sonnet (Anthropic) \and Claude Opus (Anthropic) \and ChatGPT (OpenAI)}

\date{\normalsize Hyperphysics Institute\\
\url{https://hyperphysics.com}\\
\texttt{drthomas007@protonmail.com}}

\begin{document}
\maketitle

\begin{abstract}
We prove that the Koide relation $K \equiv \Sigma m_i/(\Sigma\sqrt{m_i})^2 = 2/3$
for the three charged lepton masses follows from the adjacency spectrum of
the colour cage base graph $K_3$ (the equilateral triangle $\{V_1,V_2,V_3\}$
of the tetrahedral cage), conditional on three propositions whose epistemic
status is made explicit via a Layer~A/B/C decomposition.
(P1)~The ZBW Hamiltonian $\hat{H} = \hbar\omega_0 A_{K_3}$ is derived from
C3 cage symmetry and SSV hopping (Layer~A).
(P2)~$m_i \propto |\psi_i|^2$ is derived from the CPP DI-bit visit rate
(Layer~A).
(P3)~Equal eigenstate occupation follows from thermal equilibration
with the DP~Sea in the high-temperature limit
($kT_P/\hbar\omega_0 \sim 10^{20}$); the Gibbs formalism is
standard, but the system-bath coupling model is imported from
open quantum systems theory (Layer~B).
Under P1--P3, the eigenvalue ratio $\lambda_{\max}/|\lambda_{\min}| = 2$
forces the Lorentz modulation depth $\rho = \sqrt{2}$, from which
$K = (1+\rho^2/2)/3 = 2/3$ follows exactly (Layer~C).
At finite temperature, departures from $K = 2/3$ are of order
$\hbar\omega_0/kT_P \sim 10^{-20}$, nine orders of magnitude below
the 11~ppm experimental precision.
The result is specific to $K_3$: only the three-colour cage ($N=3$) gives
$K = 2/3$.  The theorem applies to \emph{leptons only}; quarks carry
additional strong-sector mass contributions that systematically break
Koide (deviations of 10--27\% observed).
The Koide phase $\theta$ is not determined by this theorem;
its derivation requires the electroweak sector (OP-SM-7d).
\end{abstract}

\textbf{keywords:} Koide formula, lepton mass ratio, K3 spectral theorem, 600-cell lattice, Conscious Point Physics, Zitterbewegung, charged lepton masses, Standard Model emergence, parameter-free derivation, cage base graph

%=====================================================================
\section{Background}
%=====================================================================

The Koide formula (Koide, 1982~\cite{koide1983}) is an empirical constraint
satisfied by the three charged lepton masses to 11~ppm:
\begin{equation}
K \equiv \frac{m_e + m_\mu + m_\tau}{(\sqrt{m_e}+\sqrt{m_\mu}+\sqrt{m_\tau})^2}
= \frac{2}{3}.
\label{eq:koide}
\end{equation}
No Standard Model derivation exists.  This paper proves $K = 2/3$ from
the spectral structure of the K3 cage base graph in CPP, conditional on
an open-system thermalisation model whose epistemic status is made
explicit via a Layer~A/B/C decomposition (\S\ref{sec:layers}).

\textbf{Note on independence:} The K$_3$ graph is not chosen to reproduce
the Koide relation.  It is identified as the colour cage base of the
600-cell tetrahedral cage in SM-1~\cite{abshier2026sm1} from geometric
considerations alone --- specifically, the three colour vertices
$\{V_1,V_2,V_3\}$ of the tetrahedral quark cage, whose C3 symmetry is
forced by the 600-cell lattice geometry.  The Koide relation emerges
as a consequence of that independently-motivated structure, not as a
motivation for it.

\textbf{Prerequisites used without proof here:}
\begin{itemize}
\item \textbf{Theorem~1} (Charge quantisation, SM-1~\cite{abshier2026sm1}):
  the cage base $\{V_1,V_2,V_3\}$ has C3 symmetry; $\delta = 1/3$
  (topological proof from completeness $+$ C3).
\item \textbf{Algebraic identity}: with the parametrisation
  $\sqrt{m_i} = A(1+\rho\cos\phi_i)$ and C3 phases $\phi_i = \theta+2\pi i/3$,
  one has $K = (1+\rho^2/2)/3$.  Hence $K = 2/3 \Leftrightarrow \rho = \sqrt{2}$.
\item \textbf{sea\_strength} $\approx 0.1780$ derived from the
  600-cell Voronoi geometry (SS-1, Theorem~6~\cite{abshier2026ss1}).
\end{itemize}

%=====================================================================
\section{The Colour Cage Base Graph}
%=====================================================================

\begin{definition}[Colour cage base graph $K_3$]
The three colour vertices $\{V_1,V_2,V_3\}$ of the tetrahedral cage,
with edges connecting every pair, form the complete graph $K_3$
(equilateral triangle).  Its adjacency matrix is:
\[
A_{K_3} = \begin{pmatrix}0&1&1\\1&0&1\\1&1&0\end{pmatrix}.
\]
\end{definition}

\begin{lemma}[$K_3$ adjacency spectrum]
\label{lem:K3spec}
The eigenvalues of $A_{K_3}$ are
$\lambda_{\max} = 2$ (multiplicity~1, bonding) and
$\lambda_{\min} = -1$ (multiplicity~2, antibonding).
The bonding eigenvector is $(1,1,1)/\sqrt{3}$.
\end{lemma}

\begin{proof}
$A_{K_3}(1,1,1)^T = 2(1,1,1)^T$.  For $v\perp(1,1,1)$:
$(A_{K_3}v)_i = \sum_{j\neq i}v_j = -v_i$, so $A_{K_3}v = -v$.
\end{proof}

\begin{remark}[$K = 2/3$ is unique to $N=3$]
For the complete graph $K_N$, the eigenvalues are $N-1$ (once) and $-1$
($N-1$ times).  Applying the same equipartition argument gives
$\rho^2 = N-1$ and
\[
K = \frac{N+1}{2N}.
\]
Only $N=3$ gives $K=2/3$.  The three-colour structure of the cage base
is the precise reason the lepton Koide ratio is $2/3$.
\end{remark}

%=====================================================================
\section{Epistemic Layer Structure}
\label{sec:layers}
%=====================================================================

Following the programme-wide convention established in SS-3~v1.4,
this paper makes the epistemic status of its assumptions explicit
by separating the argument into three layers, clearly distinguishing
what comes from CPP geometry, what is imported from standard
open-system quantum mechanics, and what follows mathematically.

\subsection{Layer A --- Geometric input (CPP primitives)}

The following inputs are derived from CPP axioms:
\begin{enumerate}[label=(A\arabic*)]
\item The 600-cell lattice is composed exclusively of tetrahedral
  cells (Axiom~A2).
\item Each tetrahedral cell has $3$ colour base vertices
  $\{V_1,V_2,V_3\}$ with C3 symmetry (SM-1 Theorem~1).
\item The adjacency matrix $A_{K_3}$ has eigenvalues $+2$ (once)
  and $-1$ (twice) --- a geometric fact of $K_3$
  (Lemma~\ref{lem:K3spec}).
\item The DP~Sea exists at temperature $T_P$ (Axiom~A1, A3).
\item The DI-bit exchange mechanism couples the ZBW orbital
  to the DP~Sea at each Absolute Moment (Axiom~A3).
\end{enumerate}
These are geometric and dynamical facts of the CPP lattice.
They are not postulates of this paper.

\subsection{Layer B --- Open-system formalism (imported structure)}
\label{sec:layerb}

The following assumptions are adopted from standard open quantum
systems theory.  They are \emph{not} derived from CPP primitives
in this paper:
\begin{enumerate}[label=(B\arabic*)]
\item \textbf{Caldeira--Leggett coupling.}  The ZBW--DP~Sea
  interaction is modelled as a system-bath coupling of
  Caldeira--Leggett form~\cite{caldeira1983}:
  $\hat{H}_{\mathrm{coup}} = \sum_i\sum_k g_k|i\rangle\langle i|
  (b_k + b_k^\dagger)$.
  This is a modelling choice consistent with DI-bit exchange
  but not uniquely determined by it.
\item \textbf{Rapid thermalisation.}  The relaxation time
  $\tau_{\mathrm{relax}} \ll \tau_{\mathrm{ZBW}}$, ensuring
  equilibration within each ZBW cycle.  No dynamical estimate
  of $\tau_{\mathrm{relax}}$ is given in this paper.
\item \textbf{Full Gibbs equilibration.}  The system-bath
  coupling produces the canonical Gibbs state
  $\rho_{\mathrm{eq}} = e^{-\hat{H}/kT}/Z$, not merely
  dephasing (decoherence without energy redistribution).
  This requires the bath spectral density to couple
  transitions between eigenstates of $\hat{H}_{\ZBW}$, which
  is stronger than diagonal coupling alone would
  guarantee~\cite{breuer2002}.
\end{enumerate}

\begin{remark}[Status of Layer~B]
\label{rem:layerb}
Assumptions B1--B3 are standard in open quantum systems theory
and are physically motivated by the CPP DP~Sea mechanism: every
CP exchanges DI-bits with the DP~Sea at each Absolute Moment
(Axiom~A3), providing a natural system-bath coupling.  However,
a rigorous derivation showing that DI-bit exchange dynamics
\emph{forces} Caldeira--Leggett coupling, rapid thermalisation,
and full Gibbs equilibration has not been given.  Such a
derivation --- deriving the operator formalism and system-bath
structure from CPP primitives --- would convert the present
conditional result into an unconditional one.  This is the same
Layer~B gap identified in SS-3 (Remark~3.1) and constitutes
the programme's central open problem, targeted by a future
paper (SS-4).
\end{remark}

\subsection{Layer C --- Mathematical result}

\emph{Given} Layers~A and~B:
\begin{quote}
The three $K_3$ eigenstates are equally occupied in the
high-temperature Gibbs limit.  The eigenvalue ratio
$\lambda_+/|\lambda_-| = 2$ then forces $\rho = \sqrt{2}$,
from which $K = 2/3$ follows exactly.
\end{quote}
This is proved in \S\ref{sec:theorem} (Theorem~\ref{thm:K3Koide}).

\medskip

The layered structure makes the logical dependency transparent.
Layer~C is mathematically certain.  Layer~A is geometrically
certain (given the 600-cell axiom).  The open question is whether
Layer~B can be derived from Layer~A, which would close the gap
between \emph{conditional} and \emph{unconditional} derivation
of the Koide formula.

%=====================================================================
\section{The Three Propositions}
%=====================================================================

\begin{proposition}[$\ZBW$ Hamiltonian --- derived from C3 symmetry and SSV hopping]
\label{post:H}
The ZBW orbital Hamiltonian restricted to the cage base is:
\[
\hat{H}_{\ZBW} = \hbar\omega_0\, A_{K_3},
\]
where $\omega_0 \equiv \mathrm{sea\_strength}\times c/r_{\mathrm{conf}}$.
This follows from C3 cage symmetry (Theorem~1), equal K3 edge lengths
$a = 1/\varphi$, and the SSV hopping amplitude
(Proposition~\ref{prop:P1_derivation}).
\end{proposition}

\begin{remark}[Adjacency matrix vs.\ Laplacian]
The graph Laplacian $L_{K_3} = 2I - A_{K_3}$ has eigenvalues $0$ (once)
and $3$ (twice).  The Koide-relevant ratio $\lambda_{\max}/|\lambda_{\min}|$
is the same under either the adjacency matrix or the Laplacian; the
adjacency matrix is the natural Hamiltonian for a nearest-neighbour hopping
system, where off-diagonal elements represent hopping amplitudes between
connected vertices.  The result $K = 2/3$ is not an artefact of this choice.
\end{remark}

\begin{proposition}[$\ZBW$ Born Rule --- derived from P1 and P3]
\label{prop:Born}
The mass contribution of generation $i$ is proportional to the squared
ZBW wavefunction amplitude at colour vertex $V_i$:
$m_i \propto |\psi_i|^2$.
This is derived from P1 and P3 via the CPP DI-bit flow rate
(Proposition~\ref{prop:Born_derivation}).
\end{proposition}

\begin{proposition}[Thermal equipartition --- from DP Sea thermalisation model (Layer~B)]
\label{post:equip}
At thermal equilibrium with the DP~Sea, all $K_3$ eigenstates are equally
occupied, giving $|c_+|^2 = 1/3$ and $|c_-|^2 = 2/3$.
Conditional on the Layer~B thermalisation model
(Proposition~\ref{prop:P3_derivation}; see \S\ref{sec:layerb}).
\end{proposition}

\begin{remark}[Physical motivation for P3]
In the high-temperature limit, each quantum state has equal occupation
probability $1/g$ where $g$ is the total number of states.
For $K_3$: $g = 3$ eigenstates, so each has probability~$1/3$.
The bonding sector contributes $1/3$; the two-dimensional antibonding
sector contributes $2/3$.  P3 is state-counting equipartition (equal
weight per eigenstate), not energy equipartition.
\end{remark}

\begin{remark}[Physical motivation for P1 and P2]
P1: The ZBW orbital hops between colour vertices via the electrostatic
SSV interaction.  For leptons (no colour charge), the amplitude is set
by the DP~Sea potential $V(r) = -\mathrm{sea\_strength}\times\hbar c/r$
at the confinement radius.

P2: Each CP processes DI-bit flows at rate proportional to $|\psi_i|^2$
(the ZBW occupation probability at vertex~$i$).  Mass = stored ZBW energy
$\propto$ DI-bit flow rate.  P2 is the lepton-sector analogue of the
Born rule (OP-QM-1).
\end{remark}

%=====================================================================
\subsection*{Derivation of the ZBW Hamiltonian from C3 Symmetry and SSV Hopping}
\label{prop:P1_derivation}
%=====================================================================

\begin{proposition}[C3 symmetry forces $\hat{H}_\ZBW = \hbar\omega_0 A_{K_3}$]
\label{prop:H_derived}
Given the C3 cage symmetry (Theorem~1) and equal $K_3$ edge lengths
$a = 1/\varphi$, the ZBW Hamiltonian takes the form
$\hat{H}_\ZBW = \hbar\omega_0 A_{K_3}$ with
$\hbar\omega_0 = \mathrm{sea\_strength}\times\hbar c/r_{\mathrm{conf}}$.
\end{proposition}

\begin{proof}
\textbf{Part A --- Uniformity (C3 symmetry).}
The rotation $\rho: V_1\!\to\!V_2\!\to\!V_3\!\to\!V_1$ is an exact isometry
(Theorem~1).  Any Hamiltonian commuting with $\rho$ has equal off-diagonal
elements $t_{12} = t_{23} = t_{31} \equiv t$ and equal on-site energies.
Setting the energy origin to zero gives:
\begin{equation}
\hat{H}_\ZBW = t\, A_{K_3}.
\end{equation}

\textbf{Part B --- Amplitude $t = \hbar\omega_0$ (SSV hopping).}
For leptons, the hopping amplitude is set by the electrostatic SSV
potential at the confinement radius:
\begin{equation}
t = \frac{\mathrm{sea\_strength}\times\hbar c}{r_{\mathrm{conf}}}
\equiv \hbar\omega_0 \approx 87.8~\text{MeV}.
\label{eq:hop_amp}
\end{equation}

\textbf{Conclusion.}
$\hat{H}_\ZBW = \hbar\omega_0 A_{K_3}$. \qed
\end{proof}

\begin{remark}[Consistency with SS-1 Theorem~6]
The hopping energy $\hbar\omega_0 = \mathrm{sea\_strength}\times\hbar c/r_{\mathrm{conf}}$
is consistent with the DP binding energy $E_{\mathrm{eDP}} = \hbar\omega_0/\varphi^2$
from SS-1 Theorem~6~\cite{abshier2026ss1}.  The $1/\varphi^2$ factor arises
from the Voronoi volume per vertex; $\hbar\omega_0$ is the total SSV energy
at $r_{\mathrm{conf}}$ before the per-vertex projection.
\end{remark}

%=====================================================================
\subsection*{Derivation of the ZBW Born Rule from P1 and P3}
\label{prop:Born_derivation}
%=====================================================================

\begin{proposition}[DI-bit flow rate gives $m_i \propto |\psi_i|^2$]
Given P1 ($\hat{H}_\ZBW = \hbar\omega_0 A_{K_3}$) and P3 (thermal
equipartition), the lepton mass $m_i \propto |\psi_i|^2$.
\end{proposition}

\begin{proof}
The ZBW orbital visits each colour vertex with frequency
$\Phi_i = \omega_{\ZBW} \times f_i$, where $f_i$ is the fraction of
Absolute Moments spent at $V_i$.  The mass energy stored at vertex $i$
per unit time is:
\begin{equation}
m_i c^2 = \hbar \Phi_i = \hbar\omega_{\ZBW} f_i.
\label{eq:mass_rate}
\end{equation}
By P1, the lepton generation state for generation $g$ is:
\begin{equation}
|\psi_g\rangle = c_+|\phi_+\rangle
+ c_-\, e^{i\cdot 2\pi g/3}|\phi_-^{(g)}\rangle,
\end{equation}
and the occupation at vertex $i$ is $f_i = |\langle i|\psi_g\rangle|^2 = |\psi_i|^2$.
Substituting: $m_i \propto |\psi_i|^2$. \qed
\end{proof}

\begin{remark}[P3 does double duty]
P3 simultaneously fixes $|c_-|^2/|c_+|^2 = 2$ (giving $\rho = \sqrt{2}$
and $K = 2/3$) and specifies the coherent generation state, making the
ZBW trajectory ergodic and giving $f_i = |\psi_i|^2$.
The theorem rests on \textbf{two} independent propositions (P1 and P3).
\end{remark}

%=====================================================================
\subsection*{Derivation of Thermal Equipartition from DP Sea Thermalisation}
\label{prop:P3_derivation}
%=====================================================================

\begin{proposition}[DP Sea thermalisation gives $|c_n|^2 = 1/3$]
\label{prop:equip_derived}
The ZBW K3 resonator couples to the DP Sea thermal bath at the Planck
temperature $T_P$.  Since $kT_P/\hbar\omega_0 \approx 10^{20} \gg 1$,
the Gibbs state reduces to the uniform mixture $|c_n|^2 = 1/3$
for each of the three K3 eigenstates, giving $|c_-|^2/|c_+|^2 = 2$.
\end{proposition}

\begin{proof}
\textbf{Step 1 (ZBW--DP Sea coupling --- Layer~B1).}
Every CP exchanges DI-bits with the DP~Sea at each Absolute Moment
(Layer~A5).  This exchange is modelled as a Caldeira--Leggett
system-bath coupling~\cite{caldeira1983} (Layer~B1, \S\ref{sec:layerb}):
\begin{equation}
\hat{H}_{\mathrm{coup}} = \sum_{i}\sum_k g_k\,|i\rangle\langle i|\,(b_k + b_k^\dagger).
\end{equation}
This coupling form is physically motivated by DI-bit exchange
but is not uniquely determined by it; see Remark~\ref{rem:layerb}
for its epistemic status.

\textbf{Step 2 (Thermalisation --- Layer~B2, B3).}
We assume rapid thermalisation: $\tau_{\mathrm{relax}} \ll
\tau_{\mathrm{ZBW}}$ (Layer~B2), yielding the canonical Gibbs
state (Layer~B3):
\begin{equation}
\rho_{\mathrm{eq}} = \frac{e^{-\hat{H}_\ZBW/kT_P}}{Z}.
\end{equation}
Full Gibbs equilibration (as opposed to dephasing alone) requires
the bath spectral density to couple transitions between energy
eigenstates, a condition that is assumed here but not derived
from CPP dynamics~\cite{breuer2002}.

\textbf{Step 3 (High-temperature limit --- Layer~C).}
\[
\frac{kT_P}{\hbar\omega_0} \approx 10^{20} \gg 1
\implies
|c_n|^2 \to \frac{1}{3} \quad\text{for each eigenstate}.
\]

\textbf{Step 4 (Counting --- Layer~C).}
$K_3$ has one bonding state and two antibonding states:
\[
|c_+|^2 = \tfrac{1}{3}, \quad
|c_-|^2 = \tfrac{2}{3}, \quad
\frac{|c_-|^2}{|c_+|^2} = 2. \qquad\square
\]
\end{proof}

\begin{remark}[Finite-temperature robustness]
\label{rem:robustness}
At finite temperature $T$, the exact Gibbs weights give the
antibonding-to-bonding occupation ratio
\begin{equation}
\frac{|c_-|^2}{|c_+|^2} = 2\,e^{3x},
\qquad x \equiv \frac{\hbar\omega_0}{kT}.
\label{eq:robustness}
\end{equation}
This follows from the energies $E_+ = 2\hbar\omega_0$ and
$E_- = -\hbar\omega_0$, giving Gibbs weights
$e^{-2\hbar\omega_0/kT}$ (bonding) and
$2\times e^{+\hbar\omega_0/kT}$ (antibonding, degeneracy~2),
so the ratio is $2e^{3\hbar\omega_0/kT}$.

For $x \ll 1$:
\begin{equation}
\frac{|c_-|^2}{|c_+|^2} = 2\bigl(1 + 3x + O(x^2)\bigr).
\end{equation}
At the Planck temperature, $x = \hbar\omega_0/kT_P \sim 10^{-20}$,
so the correction to the exact ratio of~$2$ is of order
$O(10^{-20})$ --- algebraically tiny, and nine orders of magnitude
below the $11$~ppm experimental precision of the Koide relation.
Departures from $K = 2/3$ are suppressed by powers of
$\hbar\omega_0/kT$, not exponentially; the robustness is
\emph{algebraic}, not doubly-exponential.
\end{remark}

%=====================================================================
\section{The Theorem}
\label{sec:theorem}
%=====================================================================

\begin{theorem}[K3 spectral origin of the Koide formula]
\label{thm:K3Koide}
Under Propositions~\ref{post:H}--\ref{post:equip}
(Layers~A and~B, \S\ref{sec:layers}) and C3 cage symmetry
(SM-1 Theorem~1), the three charged lepton masses satisfy
$K = 2/3$ exactly.
\end{theorem}

\begin{proof}
\textbf{Step 1 (Eigenstate structure, from P1).}
The ZBW eigenstates on $K_3$ are the bonding mode $|\phi_+\rangle$ at
energy $2\hbar\omega_0$ and two antibonding modes at energy $-\hbar\omega_0$.

\textbf{Step 2 (Occupation amplitudes, from P3).}
Thermal equipartition: $|c_+|^2 = 1/3$, $|c_-|^2 = 2/3$.

\textbf{Step 3 (Modulation depth $\rho$).}
\[
\rho^2 = \frac{|c_-|^2}{|c_+|^2} = \frac{2/3}{1/3} = 2
\quad\Rightarrow\quad \rho = \sqrt{2}.
\]

\textbf{Step 4 (Koide K, algebraic identity).}
With C3 symmetry (phases $\phi_i = \theta + 2\pi i/3$) and
$\rho = \sqrt{2}$:
\[
K = \frac{1 + \rho^2/2}{3} = \frac{1 + 1}{3} = \frac{2}{3}.\qquad\square
\]
\end{proof}

\begin{remark}[Physical driver]
The value $K = 2/3$ arises from the combination of the $2\!:\!1$
eigenstate degeneracy of the $K_3$ adjacency spectrum and the
high-temperature Gibbs weighting, which converts that degeneracy
into the occupation ratio $|c_-|^2/|c_+|^2 = 2$.  Both ingredients
are necessary: the degeneracy is Layer~A (geometry); the thermal
weighting is Layer~B (open-system model).
\end{remark}

\begin{corollary}[Common K3 origin of $\delta = 1/3$ and $K = 2/3$]
\label{cor:common}
Both charge quantisation ($\delta = 1/3$, SM-1 Theorem~1) and the
Koide formula ($K = 2/3$, Theorem~\ref{thm:K3Koide}) arise from the
same triangle $K_3$:
\begin{center}
\begin{tabular}{lll}
\toprule
Result & Value & $K_3$ structure used \\
\midrule
Charge fraction $\delta$ & $1/3$ & Combinatorial: 3 equal vertices, sum to~1 \\
Koide ratio $K$ & $2/3$ & Spectral: eigenvalue ratio $2:1$ \\
\bottomrule
\end{tabular}
\end{center}
\end{corollary}

%=====================================================================
\section{Why Quarks Do Not Satisfy Koide}
%=====================================================================

\begin{remark}[Quarks: Koide broken by strong-sector mass]
\label{rem:quarks}
The observed values are:
\[
K(d,s,b) = 0.731 \quad(9.7\%\text{ above }2/3), \qquad
K(u,c,t) = 0.849 \quad(27\%\text{ above }2/3).
\]
Quarks carry qDP chain binding energy ($\sim\!99\%$ of proton mass),
inter-cage bonding, and cage-depth scaling.  These strong-sector
contributions break the $K_3$ spectral symmetry.  Leptons carry no
colour charge and no qDP chains; their masses are purely from $K_3$
ZBW modes, and Koide holds to 11~ppm.
Theorem~\ref{thm:K3Koide} is a lepton result, not a universal one.
\end{remark}

%=====================================================================
\section{Scope and Open Problems}
%=====================================================================

\begin{center}
\renewcommand{\arraystretch}{1.2}
\begin{tabular}{p{5.5cm}lp{5.5cm}}
\toprule
Claim & Status & Layer / Notes \\
\midrule
$K_3$ eigenvalues $\{+2,-1,-1\}$ & \textbf{Proved} & Layer~A (geometry) \\
C3 phase structure & \textbf{Proved} & Layer~A (SM-1 Theorem~1) \\
$K = 2/3$ (Koide formula) & \textbf{Proved} & Layer~C (conditional on A+B) \\
$\rho = \sqrt{2}$ from $K_3$ spectrum & \textbf{Proved} & Layer~C \\
P1: $H = \hbar\omega_0 A_{K_3}$ & \textbf{Derived} & Layer~A (C3 + SSV hopping) \\
P2: $m_i \propto |\psi_i|^2$ & \textbf{Derived} & Layer~A (DI-bit visit rate) \\
P3: Thermal equipartition & \textbf{Conditional} & Layer~B (Gibbs model assumed; \S\ref{sec:layerb}) \\
Robustness: $\delta K \sim O(10^{-20})$ & \textbf{Proved} & Layer~C (Remark~\ref{rem:robustness}) \\
Phase $\theta$ & Open & OP-SM-7d \\
Scale $A$ & Calibrated to $m_e$ & One free parameter \\
Individual masses $m_\mu, m_\tau$ & Open & Need $\theta$ and $A$; see SM-4 \\
Caldeira--Leggett coupling & \textbf{Assumed} & Layer~B1 (modelling choice) \\
$\tau_{\mathrm{relax}} \ll \tau_{\mathrm{ZBW}}$ & \textbf{Assumed} & Layer~B2 (no dynamical estimate) \\
Full Gibbs equilibration & \textbf{Assumed} & Layer~B3 (vs.\ dephasing only) \\
\bottomrule
\end{tabular}
\end{center}

\begin{openproblem}[OP-SM-7d: Derivation of the Koide phase $\theta$]
\label{op:theta}
Derive $\theta = 132.7323°$ from CPP dynamics.
Theorem~\ref{thm:K3Koide} proves $K = 2/3$ but leaves $\theta$
undetermined.  A structural theorem (proved in SM-4,
Theorem~2~\cite{abshier2026sm4}) shows that no mechanism within the
K3+SSV framework can select $\theta$, because C3 symmetry leaves
the antibonding subspace degenerate.  The derivation of $\theta$
requires the electroweak sector (EW series, OP-EW-1).
\end{openproblem}

%=====================================================================
\section*{Acknowledgements}
%=====================================================================

Thomas Lee Abshier ND conceived the CPP framework and directed the
research programme.  Grok (xAI) contributed computational development
and the thermal picture for the ZBW resonator.  Claude Sonnet (Anthropic)
performed the formal derivations of P1, P2, and P3, verified the
proofs, and drafted versions~1--5.  ChatGPT (OpenAI) provided a
referee-grade review identifying the Layer~B gap in the P3 derivation,
the correct $O(\hbar\omega_0/kT_P)$ robustness scaling, and three
specific physics questions that sharpened the epistemic framing.
Claude Opus (Anthropic) proposed the Layer~A/B/C decomposition
adopted in version~6 and integrated the review corrections.

%=====================================================================
\bibliographystyle{plainnat}
\bibliography{../../bibliography/cpp_references}

\end{document}